\section{Introduction}
\label{sec:research-introduction}

Remote sensing (RS) imagery is one of the fastest available sources of situational
awareness after a disaster: satellites and aircraft can image an affected area
within hours, long before ground teams can reach it. Turning that imagery into
actionable information, however, still largely relies on task-specific computer
vision (CV) models trained separately for each disaster type and each output
(a flood-segmentation network here, a building-damage classifier there, a
burned-area detector elsewhere). Vision-language models (VLMs) promise a more
general alternative: a single model that can be instructed, in natural
language, to detect, segment, classify, count, and reason about damage across
disaster types and imagery sources. Whether current VLMs actually deliver on
that promise for RS disaster imagery -- and under which conditions they fail
-- is still poorly characterized.

The most direct and recent basis for this proposal is \emph{DisasterM3}
\fcite{wang2025disasterm3}, which benchmarks 14 general-purpose and RS VLMs
and fine-tunes four of them on a purpose-built, multi-hazard, multi-sensor
instruction dataset. DisasterM3 is, to date, the most comprehensive attempt
to systematically evaluate VLM capability on disaster-related RS vision
tasks. Yet a benchmark of this kind is necessarily bounded by the data and
axes of analysis its authors chose to include, and several relevant efforts
sit outside that boundary: work that appeared in parallel with DisasterM3,
work that precedes it and shaped it, and a wider landscape of
disaster-specific RS datasets that our own survey has catalogued
(Section~\ref{sec:related-work} and Section~\ref{sec:dataset-landscape}).
One gap is particularly
illustrative and motivates the direction of this proposal: DisasterM3 mixes
optical imagery from Maxar with SAR imagery from Capella Space and Umbra, and
reports that VLMs struggle with a ``cross-sensor gap'' -- but it does not
break performance down \emph{by imagery provider or sensor}, so it is not
possible to tell from the paper alone whether a model that looks strong
overall is actually failing systematically on, say, Sentinel-1 SAR or
low-resolution aerial imagery while succeeding on Maxar VHR optical. Our
accompanying dataset survey (\texttt{research\_proposal/dataset\_survey\_noora/}) records
imagery provider and sensor as first-class fields for every dataset it
catalogues precisely because this axis keeps recurring as a plausible source
of hidden performance variance, and it is a natural angle this proposal can
add on top of DisasterM3's largely provider-agnostic evaluation.

This proposal therefore asks: \emph{how well do current VLMs actually perform
on the core CV tasks of disaster response -- detection, segmentation, and
classification of flooded areas, burned areas, and damaged roads and
infrastructure (including buildings and debris) -- once evaluation is
extended across the fuller landscape of disaster RS datasets and stratified
by imagery provider and sensor, and can a consolidated benchmark or a single
fine-tuned model close the gaps that are found?} Section~\ref{sec:related-work}
positions DisasterM3 against the prior and concurrent work most relevant to
this question; Section~\ref{sec:dataset-landscape} summarizes what our
dataset survey adds to that picture; and Section~\ref{sec:research-plan}
lays out the concrete research questions and approach.
